% s-174-p174-m002; original PDF 174, 175
\Aphorism{s-174-p174-m002}{260}
Wir persönlichen Menschen leiden darunter, dass man uns im socialen Leben \Sperr{zu sehr} als Personen nimmt, als Träger bestimmter Namen und Berufe, Vertreter bestimmter Stände und Klassen, und was es noch sonst für Merkmale geben mag, um Menschen zu fixiren und abzugrenzen. Schrecklich, mit so einem Kometenschweif von Bestimmungsstücken und Erkennungsmarken durchs Leben zu wandeln, beachtet, beredet, beschwatzt, mit Fragen und Antworten und Rechten und Pflichten jeden Augenblick behelligt! „Es ist als hätte Niemand nichts zu treiben und nichts zu schaffen, als auf des Nachbarn Schritt und Tritt zu gaffen“: das ist die Seele und Moral unserer erhabenen Klatschcultur, die trotz ihrer Millionen civilisirter Menschlein gerade uns noch braucht und gerade auf uns um keinen Preis verzichten kann! — Irgendwo leben, wo man weder Nachbarn hat noch Nachbar ist; unter Völkern fremder Zunge, unbekannt, ungebunden, nichts begehrend und nicht begehrt: seinen eigenen Namen vergessen dürfen, weil Niemand danach fragt. Und zum Schluss unbemerkt verschwinden, allenfalls fünfzig Jahre später in einer Gletscherspalte zum Vorschein kommen, wohlerhalten, lächelnden Mundes — aber ohne Legitimationspapiere!
\par

